\documentclass[conference]{IEEEtran}
\usepackage[T1]{fontenc}
\usepackage{times}
\usepackage[numbers]{natbib}
\usepackage{multicol}
\usepackage[bookmarks=true]{hyperref}
\usepackage{amsmath,graphicx,booktabs,url}
\ifdefined\AnonymousSubmission
\hypersetup{pdfauthor={},pdfsubject={},pdfkeywords={}}
\else
\hypersetup{pdfauthor={Kingston Kuan and Brandon Ong},pdfsubject={},pdfkeywords={}}
\fi
\title{Beyond Dodge Rate: Testing Humanoid Dodgeball Under Camera Failures}
\ifdefined\AnonymousSubmission
\author{Author Names Omitted for Anonymous Review}
\else
\author{\IEEEauthorblockN{Kingston Kuan \quad Brandon Ong}
\IEEEauthorblockA{Hebbian Robotics\\[0.4em]
\url{https://hebbian-robotics.github.io/beyond-dodge-rate/}}}
\fi
\begin{document}
\maketitle
\begin{abstract}
A humanoid robot can dodge a ball and still leave its playing area. We study
this gap under depth-camera failures, measuring contact, balance, and area
compliance throughout each four-second trial. Two saved G1 policies, $\pi_1$ and $\pi_2$, use the
same training method and recorded settings but respond differently to stale
and blank images across 192 paired test seeds per policy. For $\pi_1$, all 122 area exits under stale images occur after the throw resolves.
In a separate fixed paired recovery test with 64 seeds, restoring fresh images
at two seconds increases success from 50.0\% to 95.3\%, while dodge rate remains 100\%.
Maximum excursion also decreases without an area cutoff. The same intervention
does not improve overall success for $\pi_2$. Shared
ball trajectories preserve the $\pi_1$ recovery benefit. These results
support evaluating motion after an encounter and testing camera handling on the
policy in use. They concern two fixed models in simulation and do not establish
a universal camera fallback or hardware safety.
\end{abstract}

\section{Introduction}
Dodge rate measures whether a robot avoids an incoming ball. It does not show
whether the robot stays upright and inside its playing area afterward.
A camera failure can therefore leave dodge rate nearly unchanged while causing
failure after the encounter. We measure all three outcomes and compare camera
inputs within two fixed policies trained with the same method and recorded settings.

Our central test restores fresh images after the simulator resolves the throw.
Matched trials share the same physical throw and camera history up to restoration,
so the comparison measures a change in subsequent motion. A control gives both
policies identical recorded throws; distance analyses check dependence on the
area requirement. Building on humanoid dodgeball~\cite{pacman,vibe} and
camera-fault evaluation~\cite{blackout}, this study tests post-encounter area
compliance and visual recovery without a new controller or training objective.

\begin{figure*}[t]
\centering\includegraphics[width=\textwidth]{generated/simulation-snapshots.png}
\caption{What dodge rate misses. Exact replays of the first test seed (3000)
for $\pi_1$, with matched views and colors changed for visibility. Top: just
before the middle robot is hit. Bottom: the four-second endpoint. The late
camera freeze allows a dodge but the robot leaves the 1 m area. The circle
marks the boundary; orange shows the pelvis path. This example illustrates
failure timing, not average performance.}
\label{fig:simulation}
\end{figure*}

\section{Methods}
\subsection{Policies and simulator}
We use Russi's G1 Depth Dodge code~\cite{depthcode}, based on ViBe/ORCS~\cite{vibe},
with LeRobot G1 deployment assets~\cite{lerobot}. The controller combines frozen
SONIC~\cite{sonic} and Theia~\cite{theia} bases with depth adaptation trained by
PPO~\cite{ppo} through RSL-RL in mjlab/MuJoCo Warp.
$\pi_1$ is the released iteration-4000 model~\cite{depthrelease}; $\pi_2$ was
trained separately by our team with the same method, architecture, and objective.
Its saved configuration matches the released reference: seed 42, 4,000 updates,
512 environments, and about 49.15 million transitions, distributed across eight
GPUs in our run. The weights differ. Initialization, random-number sequences,
and resume histories are incomplete, so these models do not constitute a
controlled comparison of training methods or repeated training runs.

Evaluation fixes the weights and native gains, torque limits, depth noise,
active-ball collisions, and fall checks. MuJoCo 3.3.7 and ONNX Runtime 1.23.2 run
on CPU, with 50 Hz control and 25 Hz images. We correct an inactive-ball bug by
disabling collisions while the ball is parked below the floor and restoring
them at launch. Twelve exact replays found no inactive-ball robot contacts or
camera visibility. All reported trials use this correction.

\subsection{Task, inputs, and matching}
Each trial lasts four seconds unless the robot falls. One ball launches at one
second. The native generator chooses among six body targets, adds jitter and
$\pm30^\circ$ direction variation, and uses a nominal flight time of 0.55--0.62 s.
Healthy depth supplies normal images; stale depth repeats the last image before
the interruption; blank depth supplies zeros. Policy history is never reset.
Fault onset is known; there is no learned failure detector.

Within each policy and seed, we require identical prelaunch robot positions and
velocities and matching launch time, target, ball position, velocity, and flight
time. We also verify camera schedules; missing or mismatched trials fail analysis.
The policies can move differently before launch and change the simulator's aim,
so equal seed numbers do not ensure equal throws across policies. Comparisons
therefore concern camera inputs within each policy, not a policy ranking.
The input-handling test fixes 192 seeds (7300--7491) per policy before observing
outcomes. Interruption starts 160 ms after launch and lasts to the endpoint.

\subsection{Outcomes and uncertainty}
Success requires dodging, staying upright, and keeping the pelvis within 1 m of
its starting position throughout the trial. We report each outcome separately.
The radius is an illustrative task requirement, not a court rule or injury
threshold. A fall uses the native checks: pelvis height below 0.45 m or upright-axis
cosine below 0.4 (about $66.4^\circ$ tilt). Aborted throws remain failures in the
total trial count. Paired 95\% intervals use 10,000 bootstrap resamples of whole
evaluation seeds, retaining their conditions together, with seed 20261014.
Intervals are per comparison, without multiplicity adjustment; per-condition
intervals are Wilson intervals. Evaluation seeds are not training replications.

\subsection{Post-encounter motion and restoration}
For all 1,152 input-handling trials, a post-hoc reanalysis locates the first
sampled area exit relative to native throw resolution: contact for hits and
native termination for misses. This is an accounting event, not closest approach.
Exits strictly afterward are late; ties count as by-resolution exits. Samples
are 20 ms apart, without interpolation; seeds and outcomes remain unchanged.

The recovery test uses five conditions: healthy, continuous stale, stale with
restoration, continuous blank, and blank with restoration. Faults start at 1.16 s;
fresh images resume at 2 s without a history reset. After a separate exploratory
16-seed pilot per policy (8000--8015), we fixed 64 new seeds (9000--9063) and the
primary contrast: restored versus continuous stale input for $\pi_1$. The other
policy and blank-input contrasts are secondary. Pilot outcomes are excluded.
All conditions retain the four-second horizon; there is no outcome-dependent
stopping. Matching checks verify identical physical traces before 2 s in each
restoration pair. All stale-input throws resolve by restoration within numerical
tolerance, so the intervention tests post-resolution motion at a fixed time.

\subsection{Shared-throw control}
An exploratory control replays 64 healthy-condition ball trajectories, taking
32 donors from each policy without outcome filtering. Both policies face the
same initial ball positions and velocities in all five camera conditions
(seeds 10000--10063; 640 trials). Native collision settings and accounting remain
unchanged; robot states can still differ. Twenty separate preflight trials and
two exact healthy-condition replays check implementation. We resample common-throw
seeds with both policies and all conditions together. This control is not pooled
into the fixed recovery estimate; it reuses recorded trajectories and does not
add independent training runs. The supplement gives donor selection and all rows.

\section{Results}
\subsection{Dodges can hide late area exits}
Under stale input, all 122 $\pi_1$ area exits occur after throw resolution;
121 follow a dodge. All 11 exits under blank input are also late. For $\pi_2$,
stale input gives seven exits by resolution and three afterward; both blank-input
exits are late. These are descriptive timings, not explanations of the weights.

\subsection{Fresh images change post-encounter motion}
\input{generated/camera-recovery-effects.tex}
Both $\pi_1$ stale conditions dodge 64/64 throws; both $\pi_2$ stale conditions
dodge 54/64. No recovery condition has falls or aborted throws. The intervention
therefore changes area compliance without rescuing the initial dodge.
Table~\ref{tab:recovery} includes all five conditions, including the smaller or
absent effects under blank input. Figure~\ref{fig:recovery} shows when exits occur.
For this $\pi_1$ and throw distribution, restored visual input reduces later
area failures; it does not establish a neural mechanism or a universal fallback.

\begin{table}[t]
\caption{Fixed recovery test: 64 seeds per policy and condition. Fresh images
resume at 2 s. Falls and aborted throws are zero. Compare conditions within
each policy.}
\label{tab:recovery}\centering\footnotesize
\input{generated/camera-recovery-table.tex}
\end{table}
\begin{figure*}[t]
\centering\includegraphics[width=\textwidth]{generated/camera-recovery.png}
\caption{All 640 fixed recovery trials. Top: successes out of 64 with Wilson
95\% intervals; paired differences are reported in the text. Bottom: fraction
of seeds that have not yet left the area, including all trials. Curves drop at
first exit even if the robot returns. Dotted: interruption at 1.16 s. Dashed:
restoration at 2 s. Bottom panels show healthy and stale conditions; the top
panels and Table~\ref{tab:recovery} include all five. Throws can differ across policies.}
\label{fig:recovery}
\end{figure*}

\subsection{Scoring choices and continuous distance}
A post-hoc check uses six radii (0.5--2 m) and four windows (2--4 s). Recovery
gains accumulate late and shrink at wider boundaries: widening from 1 to 1.5 m
at four seconds reduces the $\pi_1$ gain from 29/64 to 2/64 in the fixed cohort
and from 37/64 to 1/64 in the control.
\input{generated/recovery-motion-effects.tex}

\subsection{Stale versus blank input}
\input{generated/comparison-plain.tex}
For $\pi_1$, blank input improves area compliance from 70/192 to 181/192 but
reduces dodges from 191/192 to 182/192. For $\pi_2$, area compliance improves
from 182/192 to 190/192 while dodges fall from 165/192 to 101/192. With no falls,
these outcomes explain the opposite success effects (Table~\ref{tab:policies}),
but not their cause inside the learned policies.
\begin{table}[t]
\caption{Input-handling test: 192 seeds per policy and condition. Success
requires dodging, staying upright, and staying inside the 1 m area. All trials count.}
\label{tab:policies}\centering\small
\input{generated/policy-comparison-table.tex}
\end{table}

\subsection{Effects under shared throws}
\input{generated/common-throw-effects.tex}
The $\pi_2$ interval reaches zero, so this control alone does not establish a
harmful blank-input effect for that policy. Restoring fresh images also improves
$\pi_1$ stale-input success from 23/64 to 60/64, with 60/64 dodges in both conditions.
$\pi_2$ changes from 46/64 to 45/64, with 52/64 dodges in both. No falls or aborted
throws occur. The response difference therefore persists under identical ball
trajectories, although robot states, weights, and incomplete training histories
still differ.

\section{Related Work and Limitations}
ViBe adapts motion trackers with visual feedback~\cite{vibe}; our implementation's
released training already includes image corruption and repeated frames~\cite{depthrelease}.
PAC-MAN evaluates contact and falls with control-barrier guidance and returns
the robot to its station between throws~\cite{pacman}. Our area constraint applies
throughout a single throw. Blackout-versus-freeze work evaluates physical camera-fault
consequences in manipulation~\cite{blackout}. We extend this evaluation question
to humanoid evasion and post-encounter area compliance. VB-Com switches between
vision and blind policies~\cite{vbcom}; CAP trains robustness to degraded depth~\cite{cap}.
Our tests keep each model fixed; blank images do not make a visual model a
separately trained blind policy.

OopsieVerse separates completion from damage-free completion~\cite{oopsieverse};
our area requirement is a task constraint, not a damage measure. LIBERO-MAX uses
matched rollouts with execution-time changes~\cite{liberomax}; we match camera
conditions and verify the physical throw. Reinforcement-learning reproducibility
work documents training variability~\cite{rlrepro}, which our two models do not isolate.

The evidence covers two fixed policies, one simulated robot and starting posture,
the native throw generator, and one recorded trajectory pool. It does not establish
hardware safety, human/team interaction, a universal fallback, or training-seed
reliability. Pelvis containment does not ensure whole-body containment. Trial
archives and separately trained weights remain unreleased, limiting replication;
code, configurations, and provenance are provided through the project website.

\section{Conclusion}
A successful dodge can hide a later failure to stay inside the playing area.
For $\pi_1$, restoring fresh images after throw resolution raises success from
32/64 to 61/64 without changing the dodge count and reduces maximum excursion
without an area cutoff. $\pi_2$ responds differently. These simulation findings
support measuring post-encounter motion alongside contact and falls. Future
work should repeat the tests across independent training runs and varied throws,
then assess transfer to physical robots.

\label{page:references}
\bibliographystyle{plainnat}
\bibliography{references}
\end{document}
